\section{Методы}\label{sec:methods}
Целевая функция измеряется с аддитивным гауссовым шумом:
\begin{equation}\label{eq:f}
f(x) = \|x\|^2 + \xi, \quad \xi \sim \mathcal{N}(0, \sigma^2).
\end{equation}
Оценка градиента строится по двум измерениям в точках $x \pm c_k \Delta_k$, где компоненты
$\Delta_k$ равны $\pm 1$ с вероятностью $1/2$:
\begin{equation}\label{eq:grad}
\hat g_k = \frac{f(x_k + c_k \Delta_k) - f(x_k - c_k \Delta_k)}{2 c_k} \Delta_k^{-1}.
\end{equation}
Шаги убывают по степенному закону $a_k = a/(k+1)^{\alpha}$, $c_k = c/(k+1)^{\gamma}$ с
$\alpha = 0.602$ и $\gamma = 0.101$.
Обновление с моментом имеет вид
\begin{equation}\label{eq:momentum}
v_{k+1} = \beta v_k + \hat g_k, \quad x_{k+1} = x_k - a_k v_{k+1}.
\end{equation}
При $\beta = 0$ получается классический SPSA.
Реализация формул (\ref{eq:f})--(\ref{eq:momentum}) находится в файле \texttt{src/model.py}.
